\documentclass[11pt,a4paper]{article}
\usepackage[margin=23mm]{geometry}
\usepackage{amsmath,amssymb,amsthm}
\setlength{\emergencystretch}{2em}
\newtheorem{lemma}{Lemma}
\newtheorem{theorem}[lemma]{Theorem}
\newcommand{\gen}[1]{\langle#1\rangle}
\newcommand{\Alt}{\operatorname{Alt}}
\title{Excluding the next three first shared labels}
\author{Human Compression project: L044 continuation}
\date{}
\begin{document}\maketitle
Retain the full critical six-point-head hypotheses and all complete
predecessor proofs. Thus $n=2r\geq18$, actual ambient IP holds,
and the critical tail is the natural alternating group
$R=G_{\geq3}=\Alt((\Omega\setminus D)\cup\{6\})$.
Write
\[
 \tau=\rho_3=(6\ b)\delta,\qquad
 \sigma=\rho_4=(b\ x)\epsilon,\qquad W=D\cup\{b,x\}.
\]
The accepted pure-overlap argument makes $|W|=8$ and makes
$\delta,\epsilon$ odd involutions fixing $W$ pointwise.
The complete prefix-kernel proof supplies the strictly increasing
head orbits $C_k$, starting with $|C_3|=7$, $|C_4|=8$,
and the pure subgroup $\Alt(C_k)\leq P_k=G_{\leq k}$.
Every later unique cut has omission orbits $C_{j-1}$ and its
complement, with all shared labels smaller than $j$.

\begin{lemma}[The exterior support and its four-point orbit]\label{four}
The support of $J=\gen{\delta,\epsilon}$ has at least six points,
so $P_4$ has at least fourteen moved points.
If $\delta$ is a single transposition $(u\ v)$, the group
$J=\gen{\delta,\epsilon}$ has a four-point orbit
$U=\{u,v,u^\epsilon,v^\epsilon\}$, disjoint from $W$.
\end{lemma}
\begin{proof}
Let $m=|\delta\epsilon|$. The pure three-cycle on $\{6,b,x\}$
belongs to $P_4\cap R=\gen{\tau,\sigma}$, by actual IP and
the natural alternating groups on $W$ and on the critical tail.
A nonidentity odd-order element in this group generated by two
involutions belongs to its cyclic rotation subgroup. Since
$\tau\sigma$ has a three-cycle on $\{6,b,x\}$ and restriction
$\delta\epsilon$ outside $W$, a power producing that pure
three-cycle must have exponent divisible by $m$ and not divisible
by $3$. Hence $3\nmid m$.

The involutions $\delta$ and $\epsilon$ do not commute. If they
did, $\delta$ would centralize both $\tau$ and $\sigma$.
It also centralizes $G_{\geq5}$: those generators fix $6,b$
and commute with $\tau$. Thus it would centralize the transitive
critical tail $R$ while fixing $6$, and would be the identity,
contrary to its oddness.

This argument applies to every odd nonidentity $\delta$. Therefore
$m\geq4$ as well as $3\nmid m$. On at most five points every
odd involution is a single transposition, and the product of two
transpositions has order at most three. Thus $J$ moves at least
six points, all outside $W$. The support of $P_4$ is $W$ together
with that exterior support, proving the general assertion.
For the remaining assertion assume $\delta=(u\ v)$.

If $\epsilon$ preserved $\{u,v\}$, it would commute with
$\delta$. If it moved just one of these points outside the
pair, $\delta\epsilon$ would have a three-cycle on that point,
its image and the fixed endpoint of the pair. All other cycles
would have length at most two; its order would be divisible
by three, again impossible. Therefore $\epsilon$ moves both
points outside the pair, to distinct points $u',v'$.
On $U=\{u,v,u',v'\}$ the two generators are
$(u\ v)$ and $(u\ u')(v\ v')$, which are transitive together.
This set is invariant under both involutions and is disjoint
from $W$, proving the assertion.
\end{proof}

\begin{theorem}[Labels five and six are excluded as well]\label{next}
At a unique eligible cut of index at least five, the first shared
label cannot be five or six.
\end{theorem}
\begin{proof}
Let $g$ be the first shared label at such a cut $j$.
All shared labels are smaller than $j$, and the prefix $P_{g-1}$
fixes the other omission orbit: every earlier mark acts on the
head orbit and is unshared. In the reversed string its index
is $h=r-1-g$, the last shared index. The full interior or endpoint
support bound gives
\[
 |\operatorname{supp}(P_{g-1})|\leq
                   n-(2h+1)=2g+1.
\]
For $g\geq5$ this prefix contains $P_4$ and thus moves at least
fourteen points by Lemma~\ref{four}. For $g=5,6$ the displayed
upper bound is respectively eleven or thirteen, a contradiction.
\end{proof}

\begin{theorem}\label{exclude}
At a unique eligible cut of index $j\geq5$, the first shared
label cannot be $4$.
\end{theorem}
\begin{proof}
Suppose the first shared label is $4$, and write
$C=C_{j-1}$, $A=\Omega\setminus C$ for the omission orbits.
The prefix $P_3$ fixes $A$ pointwise, since its head labels act
on $C$ and none of its four labels is shared. Its support is
the disjoint union of $D\cup\{b\}$ and $\operatorname{supp}(\delta)$.

Reverse the entire string. The last shared index is now
$h=r-5$, and the tail past $h$ is the actual $P_3$.
The full reversed support bound, interior or endpoint as
appropriate, gives
\[
 n-|\operatorname{supp}(P_3)|\geq2h+1=2r-9.
\]
Consequently $|\operatorname{supp}(P_3)|\leq9$.
The odd nonidentity $\delta$ fixes $D\cup\{b\}$, so its support
has at least two points. Thus equality is forced, and $\delta$
is a single transposition.

First suppose $j<r-1$. The reversed crossing has interior
index $i=r-1-j\geq1$. Its earlier subgroup $G_{>j}$ fixes
$C$ pointwise, by the accepted actual prefix-orbit lemma.
The fixed set of the tail past $h$ has size exactly $2h+1$.
If that earlier subgroup is transitive on $A$, the complete
transitive-side support theorem gives $h=i+1$.
If it is intransitive on $A$, the complete intransitive-side
support theorem gives the same conclusion at equality.
All their hypotheses have the actual reversed meanings:
the second orbit $C$ is nontrivial, the earlier subgroup fixes
it, and the tail past the last shared label fixes $A$.
Hence in either case $j=5$.
But then $C=C_4$ has size eight, whereas $P_3$, supported
inside $C$, has nine moved points. This contradiction excludes
every interior cut in the statement.

It remains to consider $j=r-1$. The group $P_{r-2}$ is
transitive on $A$, fixes its complement $C$ setwise, and its
first four generators fix $A$ pointwise. Thus the restrictions
of $\rho_4,\ldots,\rho_{r-2}$ generate a transitive string action
on $A$. There are $d=r-5$ such labels, and the restriction of
$\rho_4$ is nonidentity by the first-shared hypothesis.
The endpoint $a\in A$ of the unique crossing is fixed by
every generator of index at most $r-3$.
The complete fixed-prefix/no-gap argument in the reversed
order therefore gives
\[
                         |A|\geq d+1=r-4.              \tag{1}
\]
Explicitly, a shortest path from $a$ to a point moved by
$\rho_4$ uses every intervening label $r-2,\ldots,5$;
appending its $4$-edge gives $d$ edges and $d+1$ distinct
vertices. The path exists by transitivity and nontriviality
of the last restriction. This is the minimum-length part of
the anchored-path proof, with no bound on $|A|$ assumed.

Lemma~\ref{four} now supplies the four-point $P_4$-orbit $U$
outside $C_4=W$. Both points moved by $\delta$ belong to
$C$, because $P_3$ fixes $A$. Since $P_4\leq P_{r-2}$,
the whole orbit $U$ is contained in $C$.
Let $k\geq5$ be the first index for which $C_k$ meets $U$.
The previous head orbit is disjoint from $U$, and invariance
under $P_4\leq P_k$ makes $C_k$ contain all four points.
Thus this one strict increment has size at least four.
All other increments have size at least one. Starting from
$|C_3|=7$, there are $r-5$ increments up to $C_{r-2}$,
so
\[
                          |C|\geq7+(r-5)+3=r+5.        \tag{2}
\]
Equations (1)--(2) give $n=|A|+|C|\geq2r+1$, contrary
to $n=2r$. The terminal case is excluded as well.
\end{proof}

\paragraph{Attribution and remaining scope.}
The numerical outside-support bounds and their exact interior
equality statements remain the complete acquired proofs of the
source L044 Propositions 5.20, 5.22 and 5.23, physical pages
30--33. The prefix kernels and literal omission-orbit formulas
also remain paid, with the initial parity-fibre contribution
attributed there to the root collaborator. The six-point lower
bound for the exterior support repeats the initial argument of the
acquired opposite-end proof; it does not use the second head.
The literal four-point exterior orbit and its cost in the strict
prefix growth are the present
argument. No IP is assigned to an exterior image.
Together with the independently reviewed index-three exclusions,
any shared later unique cut now has first shared label at least
seven. The other shared cuts, opposite heads, terminal singleton
and full source contracts remain separate. No whole-original,
conjecture, numerical or formalization credit is claimed.
\end{document}
